\section*{Part II: Mapping the Adjacency Space}
\addcontentsline{toc}{section}{Part II: Mapping the Adjacency Space}
\label{sec:part2_mapping_the_adjacency_space}

\subsection{The Product-Market Matrix Is the Beginning, Not the Verdict}
\label{sec:product_market_matrix_beginning}

Ansoff's matrix distinguishes market penetration, product development, market development, and diversification. It is useful because it prevents the vague claim ``we are just growing'' from hiding that the business simultaneously targets new buyers and asks them to adopt a new product. In practical analysis, however, define \textbf{market} by customer job, buying process, payer, geography, regulation, and meaningful alternatives, not only an industry code. Define \textbf{product} by delivered transformation, support obligations, technology, liability, and credible verification.

For Bayes, another short technical studio for the same working-engineer segment may be product development. Selling its existing AI teaching system to a college is market development \textbf{and service redesign}. Selling a bespoke enterprise workflow implementation is not merely a new channel for the same course; it changes the customer, success criterion, commercial risk, and operating organization. The product-market box should trigger deeper investigation, not supply a decision.

\subsection{7 Dimensions of Adjacency Distance}
\label{sec:seven_dimensions_adjacency_distance}

Use a proposed 7-dimensional distance map: \textbf{customer job, buyer and purchasing process, distribution access, production/delivery capability, talent and management, regulation and liabilities, and capital/cash timing}. Assign descriptive scores (unchanged, partially transferable, genuinely new) to each. Avoid a single weighted total before discussing critical unknowns. An initiative can be close on 6 dimensions and blocked by 1: for instance, familiar software sold into healthcare may face decisive clinical accountability and data governance obligations.

\begin{table*}[t]
\centering
\small
\caption{The 7 Dimensions of Adjacency Distance: Evaluating Execution and Governance Risk}
\label{tab:seven_dimensions_adjacency_distance}
\begin{tabularx}{\textwidth}{p{3.2cm} X p{3.6cm} p{3.8cm}}
\toprule
\textbf{Dimension} & \textbf{Core Diagnostic Question} & \textbf{Low-Distance Indicator} & \textbf{High-Distance Risk / Pitfall} \\
\midrule
1. Customer Job & Is the underlying problem identical or tightly coupled? & Familiar pain point and success metric & Different evaluation criteria and expectations \\
\addlinespace[2pt]
2. Buyer \& Procurement & Does the same person hold authorization and budget? & Established decision-maker and invoice route & Multi-stakeholder bureaucracy or committees \\
\addlinespace[2pt]
3. Distribution Access & Can existing marketing and discovery channels reach them? & Inbound demand or existing client list & Unfamiliar CAC dynamics and sales cycles \\
\addlinespace[2pt]
4. Delivery \& Quality & Does production reuse calibrated operating routines? & Reusable rubrics, tools, and workflows & Bespoke fulfillment requiring new standards \\
\addlinespace[2pt]
5. Talent \& Management & Can current staff deliver without founder intervention? & Qualified bench with available capacity & Severe key-person bottlenecks and burnout \\
\addlinespace[2pt]
6. Regulation \& Liability & Does the activity create new compliance exposure? & Unchanged legal, tax, and privacy regime & Unlicensed operations, accreditation traps \\
\addlinespace[2pt]
7. Capital \& Cash Timing & Does cash collection align with upfront delivery cost? & Prepaid fees or short receivables cycles & Working-capital drain and payment delays \\
\bottomrule
\end{tabularx}
\end{table*}

Write the actual activities to be learned. ``New market'' could mean learning 1 additional procurement process, or hiring local counsel, building multi-language support, managing tax, obtaining regulated permission, and recruiting senior country leadership. A detailed distance map makes these differences visible before the company signs contracts.

\subsection{Customer Adjacency: The Same Buyer with a New Job}
\label{sec:customer_adjacency_same_buyer}

A company can cross-sell because it already understands the customer's context and is trusted enough to be considered. Yet a buyer's trust is job-specific. A college dean who trusts an AI elective provider may not trust the same provider to redesign institutional cybersecurity. A bank that trusts a vendor's reporting dashboard may still require a separate security and risk review before awarding a sensitive automation contract.

Customer adjacency is strongest when the new job is meaningfully connected to the existing job, the buyer possesses authority and budget, and the new offer leverages information and relationships appropriately. Avoid assuming that existing customer access means unconstrained use of their data or automatic permission to market another product. Real cross-sell value comes from relevance and lower evaluation friction, not aggressive monetization of captive attention.

\subsection{Product Adjacency: Existing Capability, Different Output}
\label{sec:product_adjacency_different_output}

1 underlying capability can support many outputs. A high-quality competency assessment process could support admissions diagnostics, post-program verification, instructor calibration, employer skill audits, and curriculum evaluation. But each output changes users, interpretation, potential harm, and demand. A diagnostic used for teaching should not automatically become a hiring filter; the validity and fairness requirements differ.

The best product extension reuses tested technology, reliable content, evidence standards, or specialist routines while satisfying a separately demonstrated buyer job. The weakest extension shares little beyond a logo. Ask what proportion of design, deployment, QA, and sales must be built from scratch. More importantly, ask which existing components remain valid \textit{for the new intended use}.

\subsection{Geographic Adjacency: Same Promise, Different Conditions}
\label{sec:geographic_adjacency_different_conditions}

Geographic expansion may appear easier because the product can be distributed digitally. In practice, foreign markets have different purchasing power, taxation, consumer protection, education rules, payment friction, time zones, employment markets, language expectations, and competition. A service delivered at 8 p.m. Indian time may be accessible to some regions and impossible for others. Remote software distribution may transfer easily; in-person career services and institutional procurement usually do not.

Treat geography as an adjacency only after mapping the economic and organizational distance. A common language may reduce but not eliminate distance. Compare a small test with remote individual buyers, a local distributor, a cross-border academic partner, and a locally incorporated entity as \textit{different investment commitments}. The deeper analysis of country selection and international entry belongs to Chapters 30 and 31.

\subsection{Channel Adjacency: Using a New Route to the Same Buyer}
\label{sec:channel_adjacency_new_route}

The same product sold through direct sales, marketplaces, resellers, institutional partnerships, employers, or content-led inbound can have radically different economics. Moving into a distributor channel may broaden reach, but requires margin sharing, training, brand governance, support boundaries, and customer ownership. A new channel is not inherently additive; it may cannibalize direct buyers, trigger price conflict, or attract poorly matched customers.

Channel transfer should be tested through \textbf{incremental contribution}, not gross attributed sales. A reseller claiming 100 sales may mostly intercept customers who would otherwise buy direct. A platform may increase volume while withholding the relationship needed to establish a durable brand. A partner may improve qualified access enough to justify a significant fee. Design the evaluation to distinguish newly created demand from reassigned credit.

\subsection{Technology Adjacency: Shared Tools Do Not Prove Shared Markets}
\label{sec:technology_adjacency_shared_tools}

An organization expert in large language models might be tempted to offer AI analytics, internal search, enterprise agents, tutoring, cybersecurity tools, or medical documentation. Shared code libraries, model providers, and engineering talent reduce some development cost, but the buyer's error costs and required assurance are sharply different. The knowledge needed to build a prototype is not the same as the knowledge needed to deploy responsibly into a regulated workflow.

Map the transfer of \textit{tested capabilities}: evaluation, security, systems integration, observability, and domain expertise. Do not value an AI adjacency solely because it can be demonstrated in 1 week. A market opportunity requires paid demand, access, fulfillment, risk management, and support. Fast prototyping often hides the slow components of a business.

\subsection{Capability Adjacency: Portfolio of Skills}
\label{sec:capability_adjacency_portfolio_of_skills}

C. K. Prahalad and Gary Hamel's core-competence perspective shifts attention from organizational divisions to collective learning that supports more than 1 product. A capability such as coordinating specialized instructors, designing valid technical assessments, or handling complex regulatory customer workflows can potentially be redeployed. A list of employee résumés is not the same as a shared organizational competence: routines, interfaces, measurement, incentives, and accumulated practice make transfer possible.

Distinguish \textbf{fungible resources}, which can move easily, from \textbf{embedded capabilities}, which lose effectiveness outside the original system. A capable instructor might be excellent in a small studio but struggle teaching 300 students. An enterprise sales playbook may work in 1 industry and fail under government procurement. Estimate transfer loss and the cost of adaptation instead of assuming it away.

\subsection{Strategic Adjacency Map: Core, Near, Stretch, and Unrelated}
\label{sec:strategic_adjacency_map_four_bands}

Place candidate opportunities in 4 concentric bands:
\begin{enumerate}
    \item \textbf{Core:} Uses the same buyer and delivery model with a clearer or deeper promise.
    \item \textbf{Near adjacency:} Reuses the core capability and requires 1 or 2 new competencies.
    \item \textbf{Stretch adjacency:} Needs several new systems but could unlock meaningful future strategic options.
    \item \textbf{Unrelated:} Has little existing advantage and is best evaluated as an independent investment rather than disguised organizational synergy.
\end{enumerate}

This is a map of \textit{relative execution risk}, not a moral hierarchy. An unrelated investment can be financially rational for a mature capital allocator with cash, governance, and professional managers. It is usually a poor early-stage founder experiment if undertaken merely because the market looks fashionable. The smaller the organization, the greater the premium on coherence.

\subsection{Not All Opportunities Deserve Equal Analytical Depth}
\label{sec:not_all_opportunities_equal_depth}

Screen lightly before committing research time. A concept that violates law, the institution's ethical boundaries, or its ability to meet current obligations needs no elaborate valuation. A candidate with no identifiable buyer or realistic route to payment can be parked pending discovery. Focus detailed analysis on opportunities that satisfy basic customer, competence, and survival tests.

However, do not make screening criteria an excuse to reject all unconventional ideas. Some important adjacencies look distant precisely because they require building a new capability that becomes valuable across several future products. The proper response is to stage the experiment, cap downside, and define the new learning you expect, not to replace judgment with rigid proximity scores.
